% runtime/part_intro.tex

\chapter*{Runtime}
\phantomsection
\addcontentsline{toc}{chapter}{Runtime}

% ============================================================
\section*{What do we already know?}
% ============================================================

\begin{remark}
The previous part established that railway states do not remain static
descriptions.

They generate operational state evolution.
\end{remark}

\begin{remark}
In particular, vehicle states, balance states, comfort states, and
dynamic quantities were interpreted as states evolving relative to the
railway state.
\end{remark}

% ============================================================
\section*{What is still missing?}
% ============================================================

\begin{remark}
State evolution alone is not yet an executable engineering system.
\end{remark}

\begin{remark}
AIM requires runtime services that evaluate, reconstruct, transform,
project, cache, refine, and interpret railway states in operational
contexts.
\end{remark}

% ============================================================
\section*{What comes next?}
% ============================================================

\begin{remark}
This part introduces the runtime layer of AIM.
\end{remark}

\begin{remark}
Runtime is the operational evaluation layer that transforms stored
alignment information and railway states into usable engineering
quantities.
\end{remark}

\begin{remark}
Runtime services connect:
\begin{itemize}
\item sparse alignment semantics,
\item metric-aware geometry,
\item frame construction,
\item projection services,
\item realization-aware states,
\item dynamic quantities,
\item and vehicle-space interpretation.
\end{itemize}
\end{remark}

% ============================================================
\section*{Relation to the AIM Roadmap}
% ============================================================

\begin{remark}
Within the AIM roadmap, the present part corresponds to the transition
from dynamic state evolution toward operational runtime execution.
\end{remark}

\begin{center}
\begin{tikzpicture}[
  node distance=1.4cm,
  box/.style={
    draw,
    rounded corners,
    minimum width=4.8cm,
    minimum height=0.9cm,
    align=center
  }
]

\node[box] (dyn)
{Runtime Dynamics};

\node[box, below=of dyn] (services)
{Runtime Services};

\node[box, below=of services] (realization)
{Realization};

\draw[->, thick] (dyn) -- (services);
\draw[->, thick] (services) -- (realization);

\end{tikzpicture}
\end{center}

\begin{remark}
The progression
\[
\text{Runtime Dynamics}
\rightarrow
\text{Runtime Services}
\rightarrow
\text{Realization}
\]
marks the transition from state evolution toward operational execution
and evaluation.
\end{remark}

% ============================================================
\begin{principle}[Runtime Layer Principle]
\label{princ:runtime-layer}
Within AIM, operational quantities should be generated by runtime
operators from underlying railway states rather than stored redundantly
as independent truth.
\end{principle}

% ============================================================
\begin{summary}
Dynamics explain state evolution.

Runtime explains operational execution.

The runtime layer therefore provides the bridge between mathematical
state models and interactive, realization-aware railway engineering
systems.
\end{summary}
